%!TEX root=../../note
\subsection{Bounds}
Bounds make sense in any ordered set. The special feature of $\R$ is that every nonempty set bounded above has a supremum in $\R$; this need not hold in $\Q$.

\begin{definition}
    [Bounded Sets]
    Let $S$ be a nonempty subset of $\R$. The set $S$ is said to be 
    \begin{enumerate}
        \item bounded above: there exists $u \in \R$ such that $x \leq u$ for all $x \in S$. 
        \item bounded below: there exists $w \in \R$ such that $w \leq x$ for all $x \in S$. 
        \item bounded: it is both bounded above and below. 
        \item unbounded: not bounded. 
    \end{enumerate}
\end{definition}

If there exists an upper bound $u_0$ of $S$ such that whenever $u$ is an upper bound for $S$
we have $u_0 \leq u$, then $u_0$ is called the least upper bound, or the supremum, of $S$.
\begin{equation*}
    u_0 := \sup S
\end{equation*}

If there exists a lower bound $w_0$ of $S$ such that whenever $w$ is a lower bound for $S$
we have $w_0 \geq w$, then $w_0$ is called the greatest lower bound, or the infimum, of $S$.
\begin{equation*}
    w_0 := \inf S
\end{equation*}

\begin{remark}
    There can be only one supremum/infimum of a given set $S$ of $\R$.

    Idea: Suppose $u_1$ and $u_2$ are both suprema of $S$.

    Since $u_1$ is a supremum and $u_2$ is an upper bound, we get
    \begin{equation*}
        u_1 \leq u_2.
    \end{equation*}

    Since $u_2$ is a supremum and $u_1$ is an upper bound, we get
    \begin{equation*}
        u_2 \leq u_1.
    \end{equation*}

    Thus, $u_1 = u_2$.
\end{remark}


\begin{lemma}
    An upper bound $u$ of a nonempty set $S$ in $\R$ is the supremum of $S$ if and only if
    \begin{equation*}
        (\forall \epsilon > 0)(\exists x_\epsilon \in S)u - \epsilon < x_\epsilon.
    \end{equation*}
\end{lemma}

\begin{proof}
    First suppose $u=\sup S$. Let $\epsilon>0$. Since
    \begin{equation*}
        u-\epsilon<u,
    \end{equation*}
    the number $u-\epsilon$ cannot be an upper bound: otherwise it would be an upper bound smaller than $u$. By the definition of upper bound, there is therefore some $x_\epsilon\in S$ such that
    \begin{equation*}
        u-\epsilon<x_\epsilon\leq u.
    \end{equation*}

    Conversely, suppose $u$ is an upper bound and that for every $\epsilon>0$ there exists $x_\epsilon\in S$ with $u-\epsilon<x_\epsilon$. If $u$ were not the least upper bound, there would be another upper bound $w<u$. Choose
    \begin{equation*}
        \epsilon=u-w>0.
    \end{equation*}
    The assumed condition then yields
    \begin{equation*}
        w=u-\epsilon<x_\epsilon,
    \end{equation*}
    contradicting that $w$ is an upper bound. Hence $u=\sup S$.
\end{proof}

\begin{example}
    Consider the set $S = \set{x \in \R : 0 \leq x \leq 1}$.
    \begin{itemize}
        \item $\sup S = 1$.
        \begin{proof}
        Every $x\in S$ satisfies $x\leq1$, so $1$ is an upper bound. To see directly that no smaller number works, let $u<1$. If $u<0$, then $0\in S$ lies above $u$. If $0\leq u<1$, set $x=(u+1)/2$. Then
        \begin{equation*}
            u<x=\frac{u+1}{2}<1,
        \end{equation*}
        so $x\in S$ lies above $u$. Thus every number below $1$ fails to be an upper bound, and $\sup S=1$.
        \end{proof}
        \item $\inf S = 0$.
        \begin{proof}
            Since $x\geq0$ for all $x\in S$, $0$ is a lower bound. If $w>0$, then $0\in S$ satisfies $0<w$, so $w$ is not a lower bound. Thus no number greater than $0$ is a lower bound, and $\inf S=0$.
        \end{proof}
    \end{itemize}
\end{example}
